% !TeX program = xelatex
% UTF-8 transcription of the ten source photographs; source statements retained.
% pic01.jpg = 20261012_192026209.jpg
% pic02.jpg = 20261012_192033538.jpg
% pic03.jpg = 20261012_192040518.jpg
% pic04.jpg = 20261012_192047458.jpg
% pic05.jpg = 20261012_192053715.jpg
% pic06.jpg = 20261012_192100421.jpg
% pic07.jpg = 20261012_192106369.jpg
% pic08.jpg = 20261012_192113747.jpg
% pic09.jpg = 20261012_192119611.jpg
% pic10.jpg = 20261012_192124899.jpg
% Standalone, or input into an existing XePersian document.
\ifdefined\settextfont
  \def\STEMIStandalone{0}
\else
  \def\STEMIStandalone{1}
  \documentclass[12pt,a4paper]{article}
  \usepackage[margin=2cm]{geometry}
  \usepackage{amsmath,amssymb}
  \usepackage{xepersian}
  \IfFontExistsTF{Amiri}{\settextfont{Amiri}}{\settextfont{FreeSerif}}
  \setlatintextfont{FreeSerif}
\fi
\providecommand{\en}[1]{\lr{#1}}
\newcommand{\STEMIpage}[1]{\clearpage
  \section*{موضوع: براونوالد ـ ۳۸ ـ منج \en{STEMI}}
  {\footnotesize تاریخ: \en{/ /}\quad صفحهٔ #1\par}}
\ifnum\STEMIStandalone=1 \begin{document}\fi

\STEMIpage{۱}
\subsection*{منج \en{STEMI} در \en{ED}}
\begin{itemize}
\item نوار قلب باید \underline{کمتر از \en{10 min}} از بدو ورود اخذ شود.
\item همهٔ بیماران باید \en{IV line} و \en{CCU} داشته باشند.
\item در مرکزی که \en{PCI} ندارد، قدم اول رد کنتراندیکاسیون‌های فیبرینولیتیک است.
\end{itemize}
بیمار با \en{STD + TWI} کاندید ریپرفیوژن اورژانسی نیست مگر آنکه شک به \en{MI} خلفی داشته باشیم.

\begin{itemize}
\item زمان \en{Door to needle}: فیبرینولیتیک \en{$\leq 30\,\mathrm{min}$}.
\item \en{PCI}: \en{$\leq 90\,\mathrm{min}$}.
\end{itemize}
\subsection*{کنتراندیکاسیون‌های مطلق فیبرینولیتیک}
\begin{itemize}
\item هرگونه \en{ICH} قبلی.
\item هرگونه ضایعهٔ ساختار عروق مغز که شناخته شده باشد (\en{AVM}).
\item هرگونه بدخیمی اینتراکرانیال (اولیه، متاستاز).
\item استروک ایسکمیک ظرف \underline{۳ ماه اخیر}.
\item شک به دایسکشن آئورت.
\item خونریزی فعال (بجز منس).
\item ترومای جدی مغز و صورت ظرف ۳ ماه اخیر.
\item جراحی مغز یا نخاع ظرف ۲ ماه اخیر.
\end{itemize}

\STEMIpage{۲}
\begin{itemize}
\item \en{Severe HTN} (که به درمان اورژانسی جواب ندهد).
\item اگر ظرف ۶ ماه اخیر استرپتوکیناز گرفته باشد نمی‌توانیم استرپتوکیناز بدهیم.
\end{itemize}
\subsection*{کنتراندیکاسیون‌های نسبی}
\begin{itemize}
\item سابقهٔ \en{HTN} کنترل‌نشده: \en{$SBP>180$}، \en{$DBP>110$}.
\item شرح حال استروک ایسکمیک ظرف ۳ ماه اخیر؛ دمانس.
\item \en{CPR} طولانی‌تر از \en{10 min}.
\item جراحی ماژور ظرف ۳ هفته.
\item خونریزی داخلی ظرف ۱ ماه اخیر.
\item بارداری.
\item پپتیک اولسر فعال.
\item درمان با آنتی‌کوآگولان خوراکی.
\end{itemize}
\subsection*{درمان جنرال}
\begin{itemize}
\item آسپرین: \en{162--325 mg} خوراکی.
\item \en{ASA}: قرص \en{non enteric coated} باید جویده شود.
\end{itemize}
\subsection*{کنترل درد}
\begin{itemize}
\item مسکن‌ها: بهترین مسکن \en{M/S} است؛ تنها کنتراندیکاسیون مطلق آن حساسیت دارویی است.
\end{itemize}
\begin{latin}
\noindent Amp M/S \underline{4--8 mg} IV stat\\
\underline{2--8 mg} q 5--15 min
\end{latin}
تا رفع درد یا \en{$BP\downarrow$} یا اختلال تنفس.

\STEMIpage{۳}
\begin{itemize}
\item مورفین به دلیل خاصیت وازودیلاتوری عروق محیطی در کنترل ادم ریه نیز مؤثر است.
\item وقتی بیمار مورفین یا نیتروگلیسرین می‌گیرد باید \en{Supine} باشد و اگر فشار افت کرد، پاها را کمی بالا بدهیم.
\item در ادم ریه نباید این پوزیشن را داد.
\item اگر اثرات تحریک واگ (وازووگال) مورفین زیاد بود: \underline{آتروپین}.
\end{itemize}
\subsection*{نیترات‌ها}
\begin{itemize}
\item بهبود فلو کرونری، افزایش ظرفیت وریدی، کاهش پره‌لود.
\item در اکثر بیماران \en{ACS} نیترات زیرزبانی معقول است.
\item تنها کنتراندیکاسیون مطلق:
  \begin{itemize}
  \item شک به \en{RVMI}.
  \item \en{$SBP<90$}؛ خصوصاً در \en{MI} اینفریور یا برادی‌کارد.
  \end{itemize}
\item حتی دوز پایین آن می‌تواند \en{$BP\downarrow$} بدهد: آتروپین.
\item می‌توان از انفوزیون مداوم نیز استفاده کرد.
\end{itemize}
\subsection*{بتابلاکرها}
\begin{itemize}
\item استفادهٔ روتین از \en{BB} وریدی دیگر توصیه نمی‌شود.
\item اگر بیمار \en{HTN} داشت گزینهٔ مناسبی است.
\end{itemize}

\STEMIpage{۴}
\subsection*{پروتکل استفاده از \en{BB} در بیمار \en{STEMI}}
\begin{enumerate}
\item رد موارد زیر:
  \begin{itemize}
  \item \en{HF} نباشد.
  \item \en{$SBP<90$}.
  \item \en{$HR<60$}.
  \item \en{AV Block}.
  \end{itemize}
\item \en{5 mg} متوپرولول وریدی در ۳ دوز بولوس به فاصلهٔ ۲–۵ دقیقه.
  \begin{itemize}
  \item بعد از هر دفعه اگر \en{$HR<60$} یا \en{$SBP<100$}، دیگر دوز بعدی را نمی‌دهیم.
  \end{itemize}
\item اگر ۱۵ دقیقه پس از آخرین دوز همودینامیک پایدار بود:
  \begin{latin}
  Tab Metoprolol tartrate 25--50 q6h for 2--3 days
  \end{latin}
\item اگر دوز بالا تحمل شد: \en{100 mg BD}.
\end{enumerate}
\subsection*{اکسیژن}
\begin{itemize}
\item در بیماران اکسیژن نمی‌دهیم مگر بیمار هایپوکسیک باشد.
\item \en{O$_2$} در بیماری که هایپوکسیک نیست مضر است و \en{$BP\uparrow$} و مقاومت عروق سیستمیک ایجاد می‌کند.
\item \underline{\en{$SpO_2<90\%$}}: باید اکسیژن بدهیم.
\end{itemize}

\STEMIpage{۵}
\subsection*{انتخاب روش ریپرفیوژن}
بیمار \en{STEMI} کاندید ریپرفیوژن:
\begin{itemize}
\item مراجعهٔ اولیه به مرکز \en{PCI}: \en{Primary PCI}، \en{FMC-device $\leq 90$ min}.
\item مراجعهٔ اولیه به مرکز غیر \en{PCI} (علامت ستاره):
  \begin{itemize}
  \item اعزام برای \en{PCI}: \en{DIDO $\leq 30$ min}، \en{FMC-device $\leq 120$ min}؛ سپس \en{Primary PCI}.
  \item فیبرینولیتیک ظرف ۳۰ دقیقه از ورود:
    \begin{itemize}
    \item عدم پاسخ مناسب یا انسداد مجدد: اعزام برای \en{PCI} و مسیر \en{Primary PCI}.
    \item سایر بیماران باید ظرف \underline{۳–۲۴ ساعت} تحت آنژیوگرافی قرار بگیرند.
    \end{itemize}
  \end{itemize}
\end{itemize}
\noindent $\ast$ بیمار با شوک کاردیوژنیک یا \en{Severe HF} باید سریع تحت درمان \en{CAG} شوند، فارغ از زمان شروع \en{MI}.

انجام \en{CAG} تا ۲–۳ ساعت پس از فیبرینولیتیک ممنوع است.

\STEMIpage{۶}
\subsection*{اندیکاسیون‌های \en{CAG} در بیمارانی که فیبرینولیتیک دریافت کرده‌اند یا اصلاً درمان ریپرفیوژن نگرفته‌اند}
\begin{enumerate}
\item شوک کاردیوژنیک یا \en{Acute Severe HF} که طی بستری ما ایجاد شده باشد: \en{Class I}.
\item یافته‌های با ریسک متوسط یا بالا در تست \en{non invasive}: \en{I}.
\item \en{MI} خودبه‌خودی یا \en{early provoked}: \en{I}.
\item شکست ریپرفیوژن یا انسداد مجدد.
\item بیمار پایدار پس از ریپرفیوژن موفق ظرف ۳–۲۴ ساعت.
\end{enumerate}
\subsection*{درمان آنتی‌کوآگولان در بیمار \en{STEMI}}
\begin{itemize}
\item جدا از درمان فیبرینولیتیک است.
\item اگر بیمار بخواهد فیبرینولیتیک بگیرد:
\end{itemize}
\begin{latin}
\noindent UFH 60 unit/kg up to 4000 u IV stat and\\
UFH 12 u/kg/h up to 1000 u/h IV inf for 48 h\\
Goal APTT $\Rightarrow$ \fbox{1.5--2 ULN (50--70 s)}
\end{latin}
\begin{itemize}
\item بیمارانی که درمان فارماکولوژیک می‌شوند در طول بستری آنتی‌کوآگولان بگیرند: حداقل \underline{\en{48 h}}، حداکثر \underline{۸ روز}.
\item اگر بخواهد فیبرینولیتیک بگیرد و بدانیم بیشتر \en{48 h} می‌ماند: انوکساپارین از \en{UFH} بهتر است.
\item \en{Enoxaparin}: \en{IV bolus} + دوز \en{SC BD} می‌دهیم؛ \en{1 mg/kg}.
\end{itemize}

\STEMIpage{۷}
\begin{itemize}
\item در بیمار با سابقهٔ حساسیت به \en{HIT}: استرپتوکیناز + بیوالیرودین.
\item بیمار کاندید \en{CABG} می‌شود: \en{UFH} بهتر است.
\item بیمار می‌خواهد \en{PCI} شود: \en{UFH} یا بیوالیرودین؛ ارجح بیوا.
\end{itemize}
همهٔ بیماران \en{STEMI} با هر اپروچ باید آنتی‌پلاکت بگیرند.
\subsection*{اصول درمان آنتی‌پلاکت در \en{STEMI}}
\begin{itemize}
\item بیمار کاندید \en{PCI}:
  \begin{itemize}
  \item اگر کنتراندیکاسیون ندارد شروع آسپرین در سریع‌ترین زمان ممکن.
  \item مهارکنندهٔ \en{P$_2$Y$_{12}$} هم داده شود: کلوپیدوگرل.
  \item شروع \en{P$_2$Y$_{12}$} قبل \en{PCI} اندیکاسیون ندارد مگر آنکه هدف \en{CAG} فوری باشد.
  \end{itemize}
\item بیمار کاندید فیبرینولیتیک.
\end{itemize}
تا ۲۴ ساعت از شروع علائم می‌توان فیبرینولیتیک داد.

\underline{بهترین زمان زیر ۶ ساعت}.

\STEMIpage{۸}
\begin{itemize}
\item آسپرین: ۳۲۵ غیر انتریک کوتد جویدنی؛ سپس ۸۱ روزانه.
\item \en{P$_2$Y$_{12}$}: همهٔ \en{STEMI} با هر روش درمانی:
  \begin{itemize}
  \item ۳۰۰، ۷۵.
  \item \en{300 mg} برای بیماران زیر ۷۵ سال که فیبرینولیتیک دریافت کرده‌اند.
  \item اگر \en{Primary PCI}: \en{600 mg} قبل از آنژیوگرافی.
  \end{itemize}
\end{itemize}
\subsection*{پیشگیری ثانویهٔ \en{STEMI}}
\begin{itemize}
\item قطع مصرف سیگار.
\item جلوگیری از افسردگی: مشاورهٔ روانپزشکی.
\item \en{$LDL<55$} در تمام طول زندگی.
  \begin{itemize}
  \item کلسترول توتال تا ۴۸ ساعت پس از \en{MI} پایین می‌افتد.
  \item اگر با استاتین کنترل نشد: اضافه کردن \en{PCSK9 inh}.
  \end{itemize}
\item آنتی‌پلاکت:
  \begin{itemize}
  \item آسپرین ۸۱ در همه تا آخر عمر.
  \item حداقل ۱۲ ماه دو دارویی؛ اگر خطر خونریزی نیست، بیشتر.
  \end{itemize}
\item \en{RAAS inh}: اگر \en{HF} یا افت متوسط \en{EF} یا \en{huge RWMA}.
\item بتابلاکر: اگر کنتراندیکاسیون نیست وجود دارد یک ترایال جهت [خوانش نامطمئن: ماژوریتی] در بیماران.
\end{itemize}

\STEMIpage{۹}
\begin{itemize}
\item بیشترین تأثیر در دو تا شش هفتهٔ اول است.
\item حداقل ۲–۳ سال؛ پس اگر دلیل نیست که قطع شود ادامه تا آخر عمر.
\item نیترات‌ها: جایگاهی در پیشگیری ثانویهٔ \en{STEMI} ندارند.
\item اگر بیمار آنژین داشته باشد یا \en{HF}، باشد می‌دهیم.
\item \en{OAC}: ؟
\item \en{CCB}: تجویز روتین آن توصیه نمی‌شود.
\item در بیماران با حساسیت به \en{BB} که \en{EF} خوبی دارند: دیلتیازم، وراپامیل.
\item \en{HRT}: نباید شروع بگیرد؛ اگر از قبل بوده باید قطع بشود.
\item \en{NSAID} تجویز نشود مگر در کمترین دوز برای کمترین مدت.
\end{itemize}
\subsection*{اندیکاسیون \en{ICD}}
باید \underline{۴ هفته از \en{MI} بگذرد} و بیمار پس از ۴۸ ساعت از \en{MI}، \en{VT} یا \en{VF} نداشته باشد.
\begin{itemize}
\item \en{$EF>35\%$ $\longrightarrow$ No ICD}.
\item \en{$EF\leq 30\%$}، \en{NYHA class I} \en{$\longrightarrow$ ICD}.
\item \en{$EF\leq 35\%$}، \en{class II or III} \en{$\longrightarrow$ ICD}.
\end{itemize}

\STEMIpage{۱۰}
اگر شک داشته باشیم \en{P$_2$Y$_{12}$} تا \underline{۱۴ روز}، سپس اگر مشکلی نبود تا یک سال.

حداقل باید \underline{۶ ماه} \en{P$_2$Y$_{12}$} بدهیم.

% Printed subject/date labels and all handwritten source page numbers retained.
% The notebook emblem has no lettering.
\ifnum\STEMIStandalone=1
  \let\STEMIEnd\enddocument
\else
  \let\STEMIEnd\relax
\fi
\STEMIEnd
